\documentclass[11pt]{article}
\usepackage[margin=1in]{geometry}
\usepackage{amsmath,amssymb,amsthm,mathtools}
\usepackage{graphicx}
\usepackage[hidelinks]{hyperref}
\usepackage{microtype}

\newtheorem{theorem}{Theorem}
\newcommand{\R}{\mathbb R}
\newcommand{\dd}{\,d}

\title{A smooth one-dimensional counterexample to coordinate independence\\
in a global implicit-function growth bound}
\author{Candidate counterexample packet for arXiv:2004.04427}
\date{21 July 2026}

\begin{document}
\maketitle

\begin{center}
\fbox{\parbox{0.92\textwidth}{
\textbf{Status.} Candidate full counterexample, likely valid, pending human
review.  The growth condition in Theorem 2.1(iii) of the source is not
independent of the chosen diffeomorphisms, even for the equation $x-y=0$ on
$\R^2$ and smooth coordinate changes of $\R$.
}}
\end{center}

\section{Source question}

The source is T.~Berger and F.~Haller, \emph{On an extension of a global
implicit function theorem}, arXiv:2004.04427, published in
\emph{Comptes Rendus Math\'ematique} 360 (2022), 439--450 \cite{BH}.
Remark 2.3 on source PDF page 7 asks whether the growth bound in Theorem
2.1(iii) is independent of the diffeomorphisms $\phi$ and $\psi$.  More
precisely, if the bound holds for one pair of coordinates, must it hold for
any other diffeomorphisms $\widehat\phi,\widehat\psi$, with some continuous
weight $\widehat\omega$?  The remark explicitly leaves open whether this is
true or whether a counterexample exists.

\begin{figure}[ht]
  \centering
  \includegraphics[width=0.94\textwidth]{figures/open_problem_crop.png}
  \caption{Remark 2.3 on page 7 of the source PDF.}
\end{figure}

\section{The condition to be tested}

In the one-dimensional setting, a \emph{continuous weight} is a continuous
nondecreasing function $\omega:[0,\infty)\to(0,\infty)$ satisfying the Osgood
condition
\begin{equation}
              \int_0^\infty\frac{\dd t}{\omega(t)}=\infty. \label{eq:weight}
\end{equation}
For $F:X\times Y\to\R$, a zero set $Z$, a left inverse $S$ of $D_yF$, and
diffeomorphisms $\phi,\psi$ of the two projections, the source growth bound
is
\begin{equation}
 \bigl|D\psi(y)S(x,y)\bigr|\,
 \bigl|D_xF(x,y)(D\phi(x))^{-1}\bigr|
 \leq \omega(|\psi(y)|),\qquad (x,y)\in Z.            \label{eq:growth}
\end{equation}

\section{Proof intuition}

A global increasing diffeomorphism $h:\R\to\R$ cannot have a derivative that
is uniformly too large as a function of its output: writing $t=h(y)$, the
inverse derivative must integrate to infinite length.  This does not,
however, prevent $h'$ from having extremely high, extremely narrow spikes.

We prescribe the output-coordinate speed
$H(t)=h'(h^{-1}(t))$ to equal $1$ except on intervals of summable total
length, while making $H(n)\geq4^n$ at every positive integer.  The large
complement on which $H=1$ guarantees that $\int_0^\infty H^{-1}=\infty$, so
the resulting $h$ is a global smooth diffeomorphism.  On the other hand, any
nondecreasing function dominating $H$ must stay above $4^n$ throughout
$[n,n+1]$.  Its reciprocal is therefore integrable, so it cannot be a weight
in the sense of \eqref{eq:weight}.  Applied to $F(x,y)=x-y$, this changes an
original growth bound identically equal to $1$ into one that has no admissible
weight majorant.

\section{Counterexample}

\begin{theorem}\label{thm:counterexample}
There exist data satisfying all hypotheses of Theorem 2.1 of the source and
diffeomorphisms $\phi,\psi$ for which the growth condition
\eqref{eq:growth} holds, together with alternative smooth diffeomorphisms
$\widehat\phi,\widehat\psi$ for which it fails for every continuous weight.
Consequently the answer to the question in Remark 2.3 is negative.
\end{theorem}

\begin{proof}
Define the explicit nonnegative smooth bump
\[
 \beta(u):=
 \begin{cases}
   e\exp\!\left(-\dfrac{1}{1-16u^2}\right),&|u|<\dfrac14,\\[6pt]
   0,&|u|\geq\dfrac14.
 \end{cases}
\]
Thus $\beta\in C_c^\infty(\R)$, $0\leq\beta\leq1$, and $\beta(0)=1$.
Set
\begin{equation}
 H(t):=1+\sum_{n=1}^\infty 4^n\beta\bigl(4^n(t-n)\bigr). \label{eq:H}
\end{equation}
The sum is locally finite, hence $H$ is smooth and $H\geq1$.  The $n$th
bump is supported in
\[
 I_n:=\bigl(n-4^{-(n+1)},\,n+4^{-(n+1)}\bigr).
\]
The union $U:=\bigcup_{n\geq1}I_n$ has finite measure, since
\[
 |U|\leq\sum_{n=1}^\infty2\cdot4^{-(n+1)}=\frac16.
\]
Outside $U$ one has $H=1$.  It follows that
\begin{equation}
       \int_0^\infty\frac{\dd t}{H(t)}=\infty,        \label{eq:inverse-length}
\end{equation}
because $[0,\infty)\setminus U$ has infinite measure.  On the negative
half-line $H=1$ as well.

Now define
\[
 Q(t):=\int_0^t\frac{\dd s}{H(s)}.
\]
The function $Q$ is smooth and strictly increasing, with $Q'(t)=H(t)^{-1}>0$.
It tends to $-\infty$ as $t\to-\infty$ and, by
\eqref{eq:inverse-length}, to $+\infty$ as $t\to+\infty$.  Hence it is a
smooth diffeomorphism of $\R$ onto $\R$.  Let
\[
                         h:=Q^{-1}.
\]
The inverse-function formula gives
\begin{equation}
                         h'(y)=H(h(y)).                \label{eq:hprime}
\end{equation}

We next supply the implicit-function data.  Take
\[
 X=Y=\R,\qquad F(x,y)=x-y,\qquad
 Z=F^{-1}(0)=\{(x,x):x\in\R\}.
\]
Here the ambient Banach spaces are all $\R$.  The set $Z$ is closed and
path-connected.  Since $D_yF=-1$, the constant operator $S=-1$ is a left
inverse: $S D_yF=1$.  Choose the original diffeomorphisms
$\phi=\psi=\operatorname{id}_{\R}$.  Then the left side of
\eqref{eq:growth} is identically
\[
       |1\cdot(-1)|\,|1\cdot1|=1.
\]
Thus condition (iii) holds with the constant weight $\omega\equiv1$, which
satisfies \eqref{eq:weight}.  All hypotheses of the source theorem are
therefore satisfied.

Finally choose the alternative diffeomorphisms
\[
       \widehat\phi=\operatorname{id}_{\R},\qquad
       \widehat\psi=h.
\]
At $(x,y)\in Z$, the new left side of the growth bound is, by
\eqref{eq:hprime},
\begin{equation}
       |h'(y)(-1)|\,|1|=H(h(y)).                       \label{eq:new-bound}
\end{equation}
Suppose that some continuous weight $\widehat\omega$ made the new bound
valid.  Since $h$ maps $\R$ onto $\R$, putting $t=h(y)$ in
\eqref{eq:new-bound} would give
\begin{equation}
                       H(t)\leq\widehat\omega(|t|)
                       \qquad(t\in\R).                 \label{eq:majorant}
\end{equation}
At every positive integer, \eqref{eq:H} and $\beta(0)=1$ give
$H(n)\geq4^n$.  Since $\widehat\omega$ is nondecreasing,
\eqref{eq:majorant} implies
\[
       \widehat\omega(t)\geq\widehat\omega(n)
       \geq H(n)\geq4^n,
       \qquad n\leq t<n+1.
\]
Consequently
\[
 \int_0^\infty\frac{\dd t}{\widehat\omega(t)}
 \leq \int_0^1\frac{\dd t}{\widehat\omega(t)}
      +\sum_{n=1}^\infty4^{-n}<\infty.
\]
This contradicts the defining condition \eqref{eq:weight}.  No admissible
weight exists for the alternative coordinates, completing the
counterexample.
\end{proof}

\section{Verification report}

The construction is analytic and uses no numerical premise.  The following
points were checked explicitly:
\begin{itemize}
  \item The supports $I_n$ are disjoint and locally finite; hence the series
  in \eqref{eq:H} defines a smooth positive function.
  \item Their total length is $1/6$, so the set on which $H=1$ has infinite
  measure.  This proves the global surjectivity of $Q$, rather than merely a
  local inverse statement.
  \item The identity \eqref{eq:hprime} follows directly from differentiating
  $Q(h(y))=y$.
  \item The original data satisfy each item of Theorem 2.1, and the original
  constant weight is admissible.
  \item The contradiction uses the required monotonicity of a weight.  A
  non-monotone function could follow the narrow spikes, but such a function
  is not a weight in the source's definition.
\end{itemize}

The main human-review focus should be the quantifiers in Remark 2.3: the
original pair of diffeomorphisms satisfies condition (iii), whereas the
explicit alternative pair admits no weight at all.  Current verdict:
\textbf{full counterexample likely valid}.

\section{Scope and novelty check}

This is a counterexample to coordinate-independence of the sufficient growth
condition, not to Theorem 2.1 itself.  Indeed the underlying equation has the
global implicit function $y=x$ in every coordinate system.  The example shows
that the theorem's hypothesis can be destroyed by a smooth reparametrization
even though its conclusion remains true.

Local indexes were searched for arXiv:2004.04427, the title, ``global implicit
function'', ``growth bound'', and ``diffeomorphism independence'', with no
prior result.  A bounded web/arXiv search on 21 July 2026 used the title,
authors, DOI, exact Remark 2.3 wording, and close ``continuous weight'',
``coordinate independence'', and ``counterexample'' combinations.  It found
the source and citations but no settlement; the March 2026 arXiv text still
calls the question open.  This search was bounded, not exhaustive.

\enlargethispage{2\baselineskip}
\begin{thebibliography}{9}
\bibitem{BH}
T.~Berger and F.~Haller,
\emph{On an extension of a global implicit function theorem},
Comptes Rendus Math\'ematique 360 (2022), 439--450;
arXiv:2004.04427.  DOI: 10.5802/crmath.309.
\end{thebibliography}

\end{document}
